\begin{textsample}{POS Dim 1 – ro – Score 56.00 – ro/ro\_em\_55.txt}  \label{ex:f1_pos_142}
2.5.
A Área de Ciências da Natureza e suas Tecnologias Nos últimos anos, a defesa de \textbf{um} currículo nacional tem se intensificado principalmente no \textbf{que} se refere à garantia de direitos à educação de qualidade.
Tal discurso, encontra -se em consonância com a unificação curricular, considerando as especificações regionais.
É nesse contexto \textbf{que} a Base Nacional Comum Curricular ( BNCC ) é apresentada pelo governo brasileiro, \textbf{ancorada} na LDB e no Plano Nacional de Educação.
A proposta curricular da BNCC se constitui como \textbf{uma} \textbf{forte} política pública educacional após a reforma do Ensino Médio ( BRASIL, 2017 ).
As evoluções socioculturais e tecnológicas geram incessantes mudanças nas organizações e no nosso pensamento, \textbf{revelando} \textbf{um} novo universo no dia a dia de crianças, jovens e adultos.
Nesse contexto, \textbf{exige} -se independência para a obtenção e a seleção das informações disponíveis, bem como para a construção do conhecimento. ( GIORDAN, 2008 ).
O Ensino de Ciências da Natureza e suas Tecnologias, vem se constituindo desde meados do Século XX, sob a \textbf{influência} das mudanças de concepção do \textbf{que} é Ciência, dos avanços das pesquisas na área de desenvolvimento humano e aprendizagem, bem como do impacto da Ciência e Tecnologia no mundo \textbf{moderno} e \textbf{contemporâneo}, dentre outros fatores. ( MARTINES, 2009 ).
Segundo os Parâmetros Curriculares Nacionais ( PCNs ), o ensino de Ciências é relativamente recente e tem sido praticado de acordo com diferentes propostas educacionais, \textbf{que} se sucedem ao longo das décadas como elaborações \textbf{teóricas} e \textbf{que}, de diversas maneiras, se expressam nas salas de aula.
Onde, observamos, \textbf{que} muitas práticas, ainda \textbf{hoje}, são \textbf{baseadas} em \textbf{mera} transmissão de informações, tendo como recurso exclusivo o livro didático e sua transcrição no quadro branco ( BRASIL, 2000 ).
No Brasil, até 1961, as aulas dos componentes de Ciências da Natureza e suas Tecnologias eram ministradas apenas nas duas últimas séries do antigo curso ginasial e, apenas a partir de 1971, passou a ter caráter obrigatório nas oito séries do “ primeiro grau ”.
Ainda \textbf{que} esforços de renovação da educação estivessem sendo \textbf{preconizados} desde as décadas de 1930-40, predominava o chamado ensino tradicional, no qual cabia aos professores a transmissão de conhecimentos acumulados pela \textbf{humanidade}, por meio de aulas expositivas e diretivas, expondo o conteúdo de modo descontextualizado e sem nenhuma tentativa de diálogo com o conhecimento prévio dos alunos, cabendo a estes a reprodução fiel das informações, geralmente transmitidas e avaliadas por meio de questionários. ( LEME, 2011 ).
Esperando -se, nas avaliações orais e escritas, respostas prontas e literais dos estudantes, \textbf{que} acabavam por desestimular a tentativa de qualquer possível questionamento por estes ( LEME, 2011 ).
E, ainda, de acordo com os PCNs, no ambiente escolar o conhecimento científico era considerado \textbf{um} saber \textbf{neutro}, isento, e a verdade científica, tida como inquestionável.
A qualidade do curso era definida pela quantidade de conteúdos trabalhados ( BRASIL, 2000 ).
Entretanto, o rápido avanço das Ciências da Natureza e suas Tecnologias, seu impacto na sociedade através da tecnologia produzida a partir dos conhecimentos científicos e o papel da experimentação na Ciência \textbf{moderna}, passaram a exercer pressões por mudanças, além das questões políticas e econômicas entre os países.
As pressões destinadas para a renovação do ensino de Ciências da Natureza e suas Tecnologias foram no sentido de fazer com \textbf{que} o currículo correspondesse ao avanço do conhecimento científico e \textbf{que} incentivasse a participação ativa do estudante no processo de aprendizagem. ( GIORDAN, 2008 ).
Ainda, segundo os PCNs, os objetivos informativos deram lugar a objetivos também formativos, especialmente os relacionados às atividades práticas, \textbf{que} passaram a representar importante elemento para a compreensão ativa de diversos conceitos, mesmo \textbf{que} a sua implantação prática tenha sido difícil ( BRASIL, 2000 ).
Assim, o objetivo fundamental do ensino de Ciências da Natureza e suas Tecnologias passou o sentido de proporcionar condições para o estudante vivenciar o \textbf{que} se denominava \textbf{método} científico, ou seja, a partir de observações, levantar hipóteses, testá -las, refutá -las e abandoná -las quando fosse o caso, de forma a redescobrir os conhecimentos científicos em forma de leis e princípios, tal como fizeram os \textbf{cientistas} \textbf{que}, desde o \textbf{século} XVII, passaram a \textbf{exigir} a experimentação dentro de \textbf{uma} metodologia científica pautada pela racionalização dos procedimentos, como a indução e a dedução nas práticas de investigação de fenômenos naturais.
Visa, a alfabetização científica, desde a compreensão de conceitos e conhecimentos, da constituição social e histórica da ciência, à compreensão de questões referentes às aplicações e às implicações sociais, ambientais e éticas relativas à utilização e produção de conhecimentos científicos, à tomada de decisões frente a questões de natureza científica e tecnológica. ( SAVIANE, 2003 ; GIORDAN, 2008 ).
O estudo de cada componente curricular das Ciências da Natureza e suas Tecnologias, os quais, Biologia, Física e Química, contribui para \textbf{que} os estudantes compreendam os modos de pensar e de produzir conhecimentos próprios, os processos históricos e sociais de construção do conhecimento em cada \textbf{um} desses campos, de maneira a entenderem e construírem \textbf{posicionamentos} críticos quanto à natureza da ciência, a provisoriedade de suas \textbf{teorias} e seu papel na sociedade ( GIORDAN, 2008 ; MARCONDES, 2018 ).
A partir dos estudos de Bruner, psicólogo americano \textbf{que} se envolveu com a reforma curricular instalada nos EUA na década de 1960, ganha destaque o valor da Ciência como \textbf{uma} forma sofisticada do conhecimento humano, bem como a relevância \textbf{que} o ensino das matérias científicas deveria ter no currículo escolar.
Entretanto, diante do avanço rápido das Ciências, seria necessária \textbf{uma} \textbf{abordagem} diferente ao seu ensino.
Em vez da \textbf{exposição} dos fatos, de fenômenos e de \textbf{teorias} \textbf{cristalizadas}, Bruner defendeu a necessidade de os estudantes compreenderem o processo de descoberta científica, familiarizando -se com as metodologias das Ciências, de modo a \textbf{assimilarem} os princípios e estruturas das diversas Ciências ( BRUNER, 2001 ). \textbf{Um} aspecto \textbf{central} na \textbf{teoria} da aprendizagem de Bruner é a importância concedida ao \textbf{método} da descoberta, com base na ideia de \textbf{que} o conhecimento da estrutura das \textbf{disciplinas} \textbf{exige} a utilização das metodologias \textbf{que} suportam as várias \textbf{disciplinas} do currículo.
Com essa ideia, Bruner critica as metodologias expositivas e considera \textbf{que} a aprendizagem das Ciências ocorre através do envolvimento dos estudantes no processo de descoberta e no uso das metodologias científicas próprias de cada ciência.
A defesa do \textbf{método} da redescoberta foi realizada por importantes psicólogos, como Bruner ( 2001 ) e Piaget ( 1978 ), e levou à produção de \textbf{expressivo} material didático, especialmente relacionado à experimentação em laboratórios didáticos nas escolas, muitos deles traduzidos por especialistas para uso no Brasil e outros \textbf{que} foram produzidos por equipes locais.
Entretanto, muitas pesquisas realizadas sobre o ensino de Ciências puderam nos \textbf{revelar} o \textbf{que} muitos professores já tinham percebido, ou seja, \textbf{que} a experimentação, sem o desenvolvimento de \textbf{uma} atitude investigativa não irá garantir a aprendizagem significativa dos conhecimentos científicos com isso, novas tendências para o ensino da área começaram a emergir sob a \textbf{influência} das críticas e dos problemas sociais e ambientais, associados ao modelo desenvolvimentista, calcado na produção industrial, bem como dos problemas de saúde da população ( SAVIANI, 2003 ).
Entre as tendências para o ensino da área \textbf{que} vêm se constituindo desde os anos 1980, destacam -se as correntes “ Ciência, Tecnologia e Sociedade ” ( CTS ) e “ Ciência e Cidadania ”, \textbf{que} se \textbf{aproximam} das Ciências Humanas e Sociais Aplicadas, reforçando a \textbf{percepção} de \textbf{que} o conhecimento científico é \textbf{uma} construção humana e não \textbf{uma} verdade natural descoberta pelos \textbf{cientistas}, sendo assim, nova importância passou a ser atribuída à História e à Filosofia da Ciência no processo educacional.
Pesquisas evidenciaram \textbf{que} os estudantes possuem ideias sobre os fenômenos naturais e tecnológicos relacionados aos conceitos científicos estudados nas Ciências da Natureza e suas Tecnologias, construídas ativamente em seu meio social, independentes do ensino formal da escola. \textbf{Logo}, conhecer o \textbf{que} os estudantes já “ sabem ” sobre o conteúdo a ser ensinado e organizar situações de aprendizagem para promover mudança \textbf{conceitual}, passou a ser mais \textbf{uma} atribuição do professor desde então, A História da Ciência tem sido útil nessa proposta de ensino, pois o conhecimento das \textbf{teorias} do passado pode ajudar a compreender as concepções dos alunos do presente, além de também constituir conteúdo relevante do aprendizado.
Além disso, a preocupação de \textbf{que} o ensino das Ciências da Natureza e suas Tecnologias deve contribuir para a formação da cidadania, e esta é \textbf{uma} área privilegiada para isso, passou a integrar as novas propostas, pois, ao acreditar, durante muitos \textbf{séculos}, \textbf{que} era o centro do Universo, com a natureza à sua disposição, o ser humano, incluindo -se os riscos à sobrevivência da espécie humana, apropriou -se de seus processos, alterou seus ciclos naturais e redefiniu seus espaços, mas acabou deparando -se com \textbf{uma} \textbf{crise} ambiental global, com o aumento da temperatura \textbf{que} trouxe consigo todas as implicações, como o degelo das calotas polares, furacões, tornados, enchentes, desmoronamento de encostas, e \textbf{inúmeras} outras situações \textbf{que} colocam em risco toda vida no planeta Terra.
Bem como, a discussão das questões éticas relacionadas à associação entre desenvolvimento científico e tecnológico e interesses políticos e econômicos, levantadas no mundo \textbf{contemporâneo} diante dos problemas ambientais e do crescimento da produção, as discussões sobre os malefícios gerados pelo uso de armamento nuclear e de outras armas químicas e biológicas, de imenso potencial destrutivo e devastador, não podem ficar fora das salas de aula, principalmente do ensino de Ciências da Natureza e suas Tecnologias. ( GIORDAN, 2008 ). \textbf{Um} conhecimento maior sobre a vida e sobre sua condição singular na natureza permite ao aluno se posicionar acerca de questões polêmicas como os desmatamentos, o acúmulo de poluentes e a manipulação gênica.
Deve poder ainda perceber a vida humana, seu próprio corpo, como \textbf{um} todo dinâmico, \textbf{que} interage com o meio em sentido amplo, pois tanto a herança biológica quanto às condições culturais, sociais e afetivas refletem -se no corpo. ( MORTIMER, 2000 ). \textbf{Ademais}, passamos a conviver rotineiramente com produtos científicos e tecnológicos mais do \textbf{que} em qualquer outra época, seja para o consumo, seja para o trabalho, mas isso não significa \textbf{que} conhecemos seus processos de produção e de distribuição.
Assim, cresce a necessidade de conhecimento nesta área, para \textbf{que} \textbf{possamos} interpretar e avaliar informações, bem como participar do processo de julgamento de decisões ou divulgações científicas.
Já \textbf{que} a falta de informação científico-tecnológica pode comprometer a cidadania. \textbf{Logo}, o estudante não é só cidadão do futuro, mas já é cidadão do \textbf{hoje}, e, nesse sentido, conhecer Ciência é ampliar a possibilidade de participação para viabilizar capacidade plena de exercício da cidadania. ( BRASIL, 1998 ).
As novas propostas ou tendências para o ensino de Ciências da Natureza e suas Tecnologias, admitem a Ciência e Tecnologia como herança cultural, conhecimento e recriação da natureza.
Juntamente com a mitologia, as artes e a linguagem, a tecnologia é \textbf{um} dos aspectos fundamentais no estudo das culturas.
Atualmente, em meio à industrialização intensa e à urbanização concentradas, também potenciadas pelos conhecimentos científicos e tecnológicos.
A associação entre Ciência e Tecnologia se amplia, tornando -se mais presente no cotidiano e modificando, cada vez mais, o mundo e o próprio ser humano ( BRASIL, 1998 ).
Assim, na \textbf{contemporaneidade}, a Área de Ciências da Natureza e suas Tecnologias, antes \textbf{vista} como \textbf{um} conjunto de conhecimentos estabelecidos e prontos para serem transmitidos, passou a ser tratada como atividade humana, acentuando seu caráter experimental e processual, no qual a formação do pensamento e das atitudes dos \textbf{sujeitos} se dá nos entremeios de atividades de investigação de fenômenos, com todas as implicações políticas, econômicas e culturais relacionadas ( DELIZOICOV, 2009 ).
Ao contrário da Tecnologia, grande parte do conhecimento científico não é produzido com \textbf{uma} finalidade prática.
As Ciências Naturais, em seu conjunto, incluindo \textbf{inúmeros} ramos da Astronomia, da Biologia, da Física, da Química e das Geociências, estudam diferentes conjuntos de fenômenos naturais e geram diversas representações do mundo ao buscar compreensão sobre o Universo, o espaço, o tempo, a matéria, o ser humano, a vida, seus processos e transformações.
E ao descobrir e tentar explicar fenômenos naturais, gera conhecimento debatido e, possivelmente validado pela comunidade científica. ( BRASIL, 1998 ).
Os avanços nas discussões em torno do \textbf{que} é Ciência e os \textbf{métodos} válidos levaram os especialistas do campo do ensino de Ciências a reconhecerem \textbf{que} são as próprias \textbf{teorias} \textbf{que} “ sinalizam aos \textbf{cientistas} quais fenômenos e problemas investigar, quais \textbf{métodos} empregar. \textbf{Teorias} apresentam -se como conjunto de afirmações, hipóteses e metodologias fortemente articuladas. ” ( BRASIL, 2009, p. 24 ).
Sendo assim, as diferentes Ciências se utilizam de diferentes \textbf{métodos} de investigação e o mito do \textbf{método} científico, como \textbf{um} \textbf{método} único e igualmente significativo para todas as Ciências cai por terra.
Muitas metodologias vão sendo criadas ; às vezes, confundem -se com as próprias pesquisas.
Entretanto, alguns procedimentos são constantes na prática científica, tais como : a observação, experimentação, formulação de hipóteses ou questões de investigação, a quantificação, comparação e a busca de rigor nos resultados, e o ensino de Ciências devem desenvolver estas habilidades nos estudantes, pois elas são úteis para desenvolver o senso crítico, a reflexão e autonomia de pensamento de crianças e jovens.
Dessa maneira, as pesquisas acerca do processo de ensino e aprendizagem, os problemas sociais decorrentes do próprio crescimento da Ciência e da Tecnologia no mundo \textbf{moderno}, as reflexões da Filosofia e da História da Ciência, dentre outros fatores, levaram a várias propostas \textbf{metodológicas} para a área, geralmente reunidas sob a denominação de Construtivismo, \textbf{que} se diferencia da prática \textbf{predominante} nas escolas até os dias de \textbf{hoje}, as quais ocorrem numa \textbf{abordagem} mais instrucional, calcada nos conceitos de mudança de comportamento através do reforço. ( GIORDAN, 2008 ).
As novas tendências pressupõem \textbf{que} o aprendizado se dá pela interação professor/estudantes/conhecimento, ao se estabelecer \textbf{um} diálogo entre as ideias prévias dos estudantes e a visão científica \textbf{atual}, com a mediação do professor, entendendo \textbf{que} o estudante reelabora sua \textbf{percepção} anterior de mundo ao entrar em contato com o conhecimento científico.
As diferentes propostas reconhecem \textbf{hoje} \textbf{que} os mais variados valores humanos não são alheios ao aprendizado científico e \textbf{que} a Ciência deve ser apreendida em suas relações com a Tecnologia e com as demais questões sociais, de cidadania bem como com as ambientais. ( BRASIL, 2002 ).
Hodiernamente, estudar Ciências constitui -se em \textbf{um} modo de pensar sobre o ser humano e o mundo, analisando esta relação a partir de premissas entre novas e antigas ideias \textbf{que} sofrem \textbf{influências} do contexto social, histórico e econômico, não existindo neutralidade e objetividade absoluta, estando o conhecimento das Ciências em constante transformação.
Com isso, os modelos de ensino \textbf{embasados} em \textbf{uma} perspectiva de pura transmissão e recepção de conhecimentos, a partir da concepção indutiva da Ciência como verdade constituída e absoluta, já se encontram superados e ultrapassados, do ponto de vista epistemológico e pedagógico, pelo menos em questões referentes às proposições \textbf{teóricas}. ( CARVALHO, 2004 ; GIORDAN, 2008 ; MARTINES, 2009 ).
Sendo a Ciência \textbf{uma} construção humana, o fazer Ciência desenvolve -se em \textbf{um} processo de representação da realidade em \textbf{que} predominam símbolos e linguagem de discursos mentais do \textbf{sujeito} e discursos sociais do coletivo. \textbf{Logo}, o objetivo da aprendizagem, na perspectiva de Bruner, deve ir além de propiciar o prazer do domínio de \textbf{um} dado conteúdo, devendo sim, ser útil no presente e no futuro do indivíduo, sendo essencial estar no âmago do processo educativo, levando ao contínuo aprofundamento e à ampliação do saber, desenvolvendo atitudes frente à investigação, acreditando na possibilidade de utilizar o conhecimento para descobrir regularidades e semelhanças entre ideias e fatos ( LEME, 2011 ).
Atualmente, sabemos, por \textbf{intermédio} de vários estudos, \textbf{que} a aprendizagem é alcançada por meio da mediação dos significados, \textbf{que} surgem nas interações \textbf{que} se estabelecem entre os estudantes, o professor e o conhecimento durante as aulas e/ou a partir delas.
Nessa perspectiva, o conceito de aprender está diretamente ligado a \textbf{um} \textbf{aprendiz} \textbf{que}, por suas ações, envolve seus colegas e seu professor, busca e adquire informações em fontes variadas, ressignificando e produzindo reflexões e transformações do conhecimento apropriado no processo.
Portanto, entende -se \textbf{que} a fonte da dificuldade de aprendizagem em Ciências não \textbf{depende} da natureza do conteúdo, mas sim, de sua apresentação de modo muito formalizado ( \textbf{FREIRE}, 1998 ; BRUNER, 2001 ).
Recentemente, está em debate a aprendizagem colaborativa no ensino de Ciências da Natureza, por meio da qual são criadas oportunidades para promover a colaboração entre as equipes e não somente realizar experimentos, tendo lugar a \textbf{contextualização} de temas \textbf{que} sejam socialmente significativos para a aprendizagem dos alunos, como por exemplo, as problemáticas ambientais locais, regionais e mundiais, tanto para a \textbf{problematização} ( temas relevantes ) quanto à organização do conhecimento científico ( epistemologicamente significativos ), sendo atribuído ao professor o papel principal de coordenador das atividades.
Problematizando, o professor aproveita situações de questões \textbf{que} possam ser colocadas, sem, no entanto, fornecer respostas prontas, evitando a acomodação de seus alunos.
Utilizando -se mais do questionamento, estimulando a autonomia do aluno e promovendo, segundo Bruner, a aprendizagem por descoberta ( LEME, 2011 ).
Bruner enfatiza a aprendizagem por descoberta, porém, de \textbf{uma} maneira"dirigida", de modo \textbf{que} a exploração de alternativas não seja caótica ou cause confusão e angústia no estudante.
Se, por \textbf{um} \textbf{lado}, \textbf{um} roteiro de estudo, por exemplo, não deve ser do tipo"receita de bolo", por outro, não deve também ser totalmente desestruturado, deixando o aluno sem saber o \textbf{que} fazer.
As instruções devem ser dadas de modo a explorar alternativas \textbf{que} levem à solução do problema ou a sua descoberta.
As ideias básicas de \textbf{uma} Ciência podem ser captadas inicialmente de modo intuitivo, para depois serem formalizadas, à medida \textbf{que} o desenvolvimento cognitivo evolui, e são retomadas em atividades escolares subsequentes em \textbf{um} currículo estruturado, como o currículo em espiral proposto por Bruner ( 1978 ), \textbf{que} enuncia \textbf{que} o \textbf{aprendiz} deve ter oportunidade de \textbf{ver} o mesmo \textbf{tópico} mais de \textbf{uma} vez em distintos níveis de profundidade e modos de representação e, \textbf{que} qualquer Ciência pode ser ensinada, pelo menos nas suas formas mais \textbf{simples}, a alunos de todas as idades, \textbf{uma} vez \textbf{que} os mesmos \textbf{tópicos} serão, posteriormente, retomados e aprofundados mais tarde.
Desse modo, \textbf{uma} das frases mais \textbf{conhecidas} de Bruner confirma “ é possível ensinar qualquer assunto, de maneira honesta, a qualquer criança, em qualquer estágio de seu desenvolvimento. ” ( BRUNER, 1978 ).
Para tanto, \textbf{um} dado conteúdo de ensino deve ser concebido como conhecimento próprio para informar e orientar o juízo prático, sendo \textbf{que}, havendo tal interesse, possibilitará a comunicação entre os \textbf{sujeitos} envolvidos na aprendizagem, considerando também as condições objetivas do conhecimento a ser alcançado.
Assim, segundo GIORDAN ( 2008 ), \textbf{uma} das principais funções do ensino de Ciências da Natureza e suas Tecnologias na escola, deve ser o de permitir aos estudantes a apropriação do conhecimento científico ( MORAN, ( 2009 ).
Portanto, o foco na organização das competências e habilidades da área de Ciências da Natureza é contribuir para \textbf{que} o estudante, ao finalizar o Ensino Médio, alcance a formação integral.
Nesse sentido, alguns temas, \textbf{que}, porventura, não estejam contemplados no organizador curricular, podem ser inseridos pelo professor da escola, tendo este o conhecimento da cultura, tradições e da realidade local, corroborando a nova proposta do Ensino Médio.
Para tanto, este documento contempla as dez competências gerais propostas para o ensino médio pela BNCC, as três competências específicas da área de Ciências da Natureza, bem como as vinte e sete habilidades da área.
Assim, o estudante poderá apropriar -se das diversas linguagens, com o intuito de possibilitar a compreensão dos fenômenos naturais e sua interpretação nas aulas. \textbf{Ademais}, é necessário \textbf{que} o estudante conheça os \textbf{fundamentos} da tecnologia \textbf{atual}, já \textbf{que} ela atua diretamente em sua vida e \textbf{certamente} definirá o seu futuro profissional.

% matched lemmas: abordagem, ademais, ancorar, aprendiz, aproximar, assimilar, atual, basear, central, certamente, cientista, conceitual, conhecido, contemporaneidade, contemporâneo, contextualização, crise, cristalizar, depender, disciplina, embasar, exigir, exposição, expressivo, forte, freire, fundamento, hoje, humanidade, influência, intermédio, inúmero, lado, logo, mero, metodológico, moderno, método, neutro, percepção, posicionamento, possar, preconizar, predominante, problematização, que, revelar, simples, sujeito, século, teoria, teórico, tópico, um, ver
\end{textsample}
